%! Summary Statistics for a Quantitative Variable — Unit 1, Lesson 1.5 — Student Packet Cover Sheet
\documentclass[12pt]{article}
\usepackage{apstats-article}
\usepackage{apstats-boxes}
\usepackage{ltablex}
\keepXColumns

\begin{document}

% Full-bleed navy banner. \coverbanner MEASURES the title block and sizes the
% band to it, so a lesson title long enough to wrap grows the band rather than
% running out the bottom of it.
\coverbanner{Unit 1}{Lesson 1.5 \quad Summary Statistics for a Quantitative Variable}

\namedateperiod
\vspace{0.18in}

\boxguard[14]
\begin{learningtargetbox}
\small
% GRADUAL RELEASE: the vocabulary is taught in the guided notes, so the targets name it
% plainly. The AP Learning Objectives (UNC-1.I, UNC-1.J, UNC-1.K) stay in the teacher-facing plan.
\begin{enumerate}[leftmargin=1.5em, itemsep=5pt]
  \item \ldots{} calculate the \textbf{mean} and the \textbf{median} of a quantitative
        variable, and find the \textbf{quartiles} $Q_1$ and $Q_3$ by splitting an ordered list.
  \item \ldots{} read a \textbf{percentile} in context, and put a single number on spread ---
        the \textbf{range}, the \textbf{interquartile range}, or the \textbf{standard deviation}.
  \item \ldots{} say which summaries one unusual value can hijack, using the words
        \textbf{resistant} and \textbf{nonresistant}.
  \item \ldots{} choose which pair of summaries to report and defend the choice from the
        \textbf{shape} of the distribution.
\end{enumerate}
\end{learningtargetbox}

\vspace{0.14in}

\boxguard[14]
\begin{tocbox}
\small
\renewcommand{\arraystretch}{1.5}
\begin{tabularx}{\linewidth}{@{} c l X r @{}}
\rowcolor{sky}
\textbf{\#} & \textbf{Component} & \textbf{Description} & \textbf{Score} \\
\midrule
1 & Warm-Up & A shape from a dotplot, an average, and one ordered list cut into quarters & \blank{1.2cm} \\
2 & Guided Notes \& Practice & Nine paychecks, three delivery routes, and the summaries that
                 move when one value does --- then eleven quiz scores worked together & \blank{1.2cm} \\
% AP Practice is assigned but NOT scored: its score cell prints NA, never a \blank{}.
% Row 4 is the lesson's GRADED practice. Paper homework is the DEFAULT for this course; on the
% occasions DeltaMath is assigned instead, that replaces row 4 and its score cell is struck.
3 & AP Practice & Household water use in 15 homes: AP-style multiple choice and one
                 free-response set --- practice, not a grade & \textbf{NA} \\
4 & Homework & Commutes, summer books, homework minutes, a percentile, and one multiple-choice
                 item --- \textbf{scored, due the first class after two study halls} & \blank{1.2cm} \\
\midrule
  & \multicolumn{2}{r}{\textbf{Total}} & \blank{1.2cm} \\
\end{tabularx}
\end{tocbox}

\vspace{0.14in}

\boxguard[8]
\begin{remindbox}
\small
Nine paychecks have a mean of \$1{,}000 and a median of \$750, and \textbf{both are correct}.
Swap the owner's \$3{,}000 for \$750 and the mean falls \$250 while the median does not move:
mean, standard deviation, and range are \textbf{nonresistant}; median and IQR are
\textbf{resistant}. The shape picks the pair you report --- skewed or an outlier present, median
and IQR; roughly symmetric with no outliers, mean and standard deviation. And a range is
\textbf{one number}, maximum minus minimum --- never ``from 27 to 51.''
\end{remindbox}

\end{document}
